\chapter{Limitations and Future Work} \label{chapter:limitationsfutureworks}

This chapter discusses the limitations of the current research which stem from the various methodological decisions taken as part of the research process in Section \ref{sec: limitations}. Additionally, the plausible research pathways that can be explored in the future are discussed as part of Section \ref{sec: future work}.

\section{Limitations} \label{sec: limitations}

As discussed in Section \ref{Subsec: recruiting}, using an English-only instrument and snowball sampling meant that the sample was made up of English-speaking, largely urban, educated and relatively young users. While they would be the first to see the after-effects of a digital personal data privacy based legislation like DPDPA, it meant that the opinions of the under-represented minorities could not be understood. This study is explicitly posited as an exploratory, non-representative work, non-generalizable work. The vast linguistic differences present in India could therefore not be captured as well. For these reasons, all the results discussed throughout this study must be taken with respect to only this specific sample. Telangana also accounts for 33\% of the users in the sample, possibly skewing the results of the overall sampled population. As mentioned in Section \ref{Subsec: Data collection}, the results for all the analyses minus the Telangana sample can be found in the Appendix.

Previous literature which use the IUIPC instrument to measure the general privacy concerns of users report high mean values or ceiling effect being present \cite{gros_validity_2021, malhotra_internet_2004, 10.1145/3380960, 10.1145/3706598.3713912}. This was also noted in the pilot study that was conducted \ref{subsec: pilot analysis}. To account for this, the third item in the awareness sub-scale was reverse-worded, to ensure respondent's attention was accurate and not primed towards the higher scores. Unfortunately, this could not be achieved, and in turn also revealed the inattentiveness of some respondents while answering this specific section, which resulted in the bi-modal response distribution for the reverse-keyed item \ref{fig:IUIPC Frequency}. Additionally, the observed ceiling effect on some of the personal data items \ref{fig:Frequency sensitivity} weakened the correlations, potentially inflating the null demographic effect results. 

The homogeneous sample skewing slightly towards male, technical users meant that the isolated effects of each demographic could not be properly observed. The null result in both the Mann-Whitney U test and the regression models, therefore must be considered with this caveat in mind.

Additionally, using self-reported data obtained from the Likert scales means the possibility of recollection bias from the respondents. Though this survey was not incentivized, which reduces inattentive responses for monetary benefits, an attention check question added into the personal data items helped weed out 40 responses who failed the said check \ref{subsec: refinements}. Moreover, the quality and depth obtained from the approximately 100 open-text answers for each of the 6 questions showcase user interest in engaging with this conversation. Unfortunately, these open-text questions cannot be a replacement for actual qualitative study in terms of richness and hence the data from the said questions only offer comparatively limited insights.

Finally, the DPDPA itself is a relatively new framework under partial enforcement at the time of the survey. While this novelty helps with capturing the raw insights from the participants, it also means that the awareness figures must only be considered as a snapshot situated in the current setting \ref{tab:dpdpa-awareness}. Coupled with the political unrest in the country which was discussed in \ref{DPDPA discussion}, situational priming of the users towards institutional distrust cannot be ruled out.


\section{Future work} \label{sec: future work} 

This study primarily focuses on respondents living in India, since they are the main focus of the Act and our study. However, DPDPA is also applicable based on where the data is processed, not just who the person is. So, Non-resident Indians (NRIs) and Overseas Citizens of India (OCIs) when using Indian services can use its rights (access, correction, erasure, grievance redressal). Their views on privacy may still be shaped by the laws of the countries they currently live in, such as the GDPR and CCPA. 

Snowball sampling meant that responses from some users who currently do not reside in the country were also collected (n=39). As seen in Section \ref{subsub:demo effects}, this NRI sub-sample shows statistically significant differences in their privacy sensitivity of \textit{Medical records and history} and \textit{Passport details}, while having lesser perceived control than their domestic respondents. However, this sample is small, so these results are for descriptive purposes only. Future works should use a larger sample, additionally accounting for migration duration, regulatory environment differences, and personal and cultural adaptation. This is worth exploring, given the recent shift in data privacy \cite{edpb2024opinion} and transparency \cite{eu2024aiact} debates with respect to LLM training in the Global North, which indirectly influences user data privacy perceptions.  

Additionally, as discussed in Section \ref{sec: resignation}, this sample showcased Digital Resignation in their perspectives regarding their personal data, which primarily stemmed from the distrust in the corporations and legal hurdles. Future works can potentially run studies to explicitly understand this phenomenon and how it is showcased in the Global South contexts with even larger samples. Moreover, using validated scales for attitudes like Privacy Cynicism \cite{hoffmann_privacy_2016, lutz_data_2020} can help in understanding how this interconnectedness of corporate motives and legal barriers might influence user behavior in the Global South.  